	\vspace{12pt}
	\renewcommand{\labelenumi}{\alph{enumi}.)}
	\section{Am Tag der Prüfung}
		Jedes Mitglied, welches an einer Prüfung der Karate Abteilung der \href{https://www.sglangenfeld.de/}{SGL}, Standort Reusrath, teilnehmen möchte, sollte, neben den Grundvoraussetzungen, ebenso einige Leitlinien am Tag der Prüfung befolgen.\vspace{4pt}\newline
		Zunächst ist es wichtig, dass sich jeder Trainierende vor seiner Gürtelprüfung gut und richtig vorbereitet fühlt. Dazu gehören auch feste Absprachen bei der Partnerwahl für Kihon-Ippon-Kumite, Nage-Waza, Kumite-Ura, Bunkai, SV und Kumite (Partnerformen!). Diese werden nicht erst am Tag der Prüfung gewählt, sondern genau wie die anderen Prüfungsaspekte vorher im Dojo abgesprochen und demnach trainiert.\vspace{4pt}\newline
		Im Falle von körperlichen Schwächungen, Unwohlsein, Verletzungen, etc. muss der Trainer/Prüfer darüber \mbox{\textbf{grundsätzlich}} in Kenntnis gesetzt werden. Dieser entscheidet dann ob
		\begin{enumerate}[nosep]
			\item Der Trainierende von der Prüfung ausgeschlossen wird und an einem späteren Termin die Prüfung nachholen kann
			\item Der Trainierende an der Prüfung teilnehmen kann, jedoch auf einzelne Prüfungsaspekte verzichten muss oder diese in abgeschwächter Form zeigen soll
		\end{enumerate}
		\begin{center}
			\textbf{Alle Teilnehmer der Prüfung haben zudem ordnungsgemäß und pünktlich auf der Prüfung zu erscheinen.}
		\end{center}
		Ordnungsgemäß bedeutet dabei:
	\renewcommand{\labelenumi}{$\bullet$}
		\begin{itemize}[nosep]
			\item Gewaschener und gebügelter Dogi (Anzug)
			\item Ein weißes T-Shirt unter dem Dogi der Frauen/kein T-Shirt unter dem Dogi der Männer
			\item Zu einem Zopf zusammengebundene, lange Haare
			\item Saubere, gewaschene Füße
			\item Kurze und gepflegte Finger-, und Fußnägel (Verletzungsgefahr!)
			\item Kein Schmuck in jeglicher Form
			\item Ein gepflegtes und sauberes Äußeres im Allgemeinen
		\end{itemize}
		Werden diese Aspekte nicht beachtet, hat der Karate-Ka die Option bis zum Beginn der Prüfung alles zu richten. Andernfalls kann an der Prüfung nicht teilgenommen werden. Erscheint der Prüfling zudem verspätet, beginnt die Prüfung nichtsdestotrotz zur angesetzten Uhrzeit. Nur entschuldigte Verspätungen werden demnach in die Abfolge der Prüfung einkalkuliert.
	\vspace{4pt}\newline
		Die Prüfungen finden in einem offenen Rahmen statt, das bedeutet, das Familie, Freunde und Bekannte herzlich eingeladen sind, sich die Prüfung anzuschauen. Voraussetzung hierfür ist, dass den Prüflingen trotzdem eine ruhige Atmosphäre geboten werden muss. Fotos/Videos dürfen ebenso gemacht werden, nur laute und störende Geräusche durch elektrische Geräte sollten vermieden werden.
	\vspace{4pt}\newline
		Die Prüfungsgebühr und die Gebühr für die Prüfungsmarken (für den schwarzen Pass des \href{https://www.karate.de/}{DKV}, siehe dort Beitrags-, und Gebührenordnung) werden entweder einige Tage vor der Prüfung bezahlt oder am Tag der Prüfung bei den Prüfern. Wie dies gehandhabt wird, geben die Trainer frühzeitig bekannt. Die Abfolge der Prüfung gestaltet sich so, wie der Prüfer am selben Tag veranlasst. In aller Regel beginnt das Prüfungsprogramm der Teilnehmer mit den niedrigsten Gürtelfarben. Die Aspekte, welche im karatespezifischen Teil dieses Dokuments aufgeführt sind, werden demnach geprüft. Nachdem alle Prüflinge ihr Programm vorgeführt haben und der Prüfer die Bewertungen (\textit{diese werden grundsätzlich weder veröffentlicht noch zur persönlichen Kenntnis gebracht}) eingetragen hat, wird veröffentlicht, wer fortan den nächsten Kyu-Grad tragen wird.
	\vspace{4pt}\newline
		Besteht jemand die Prüfung nicht, kann zu jeder Zeit ein neuer Versuch gestartet werden. Entweder zu einem regulären Termin des Standorts, oder, in Absprache mit den Trainern, zu einem Termin und Ort außerhalb des Dojos (auf Lehrgängen oder Prüfungen in anderen Dojos des \href{https://www.karate.de/}{DKV}).
	\vspace{4pt}\newline
		Besteht der Karate-Ka die Prüfung, trägt er ab diesem Zeitpunkt die nächsthöhere Gürtelfarbe, bzw. reiht sich entsprechend seines Kyu-Grades links neben den höheren Kyu-Graden ein.
	\vspace{4pt}\newline
		\begin{flushright}
			gez. die Trainer und Prüfer der G\={o}j\={u}-Ry\={u} Karate Abteilung der \href{https://www.sglangenfeld.de/}{SGL}, Standort Reusrath
		\end{flushright}